\documentclass[11pt]{article}
\usepackage[margin=1in]{geometry}
\usepackage{amsmath,amssymb}
\usepackage{booktabs}
\usepackage{hyperref}
\usepackage{xcolor}
\usepackage{empheq}

\newcommand{\code}[1]{\texttt{#1}}
\newcommand{\file}[1]{\texttt{#1}}
\newcommand{\vect}[1]{\boldsymbol{#1}}

\title{The magnification-PDF emulator: an equation-level pipeline note}
\author{\texttt{gw-wl-emulator/production\_model}}
\date{\today}

\begin{document}
\maketitle

\noindent\textit{Verified in this session against \code{gw-wl-emulator} commit
\code{fcd795b} (2026-09-04), by reading the shipped
\code{production\_model/} package directly (not the exploratory
\code{gw-wl-emulator-ar} research tree) and by loading the two frozen
checkpoints to confirm their recorded hyperparameters. Numbers quoted below
come from files under \code{production\_model/evaluation/results/},
\code{production\_model/model/calibration/}, and
\code{production\_model/PROVENANCE.md} in that same commit. A later retrain
or a different checkpoint pairing will not automatically match this note.}

\begin{abstract}
This note documents, equation by equation, what the shipped magnification-PDF
emulator (\code{production\_model/}) actually computes, and why, for a reader
who is not an ML specialist. The paper's Sec.~III (``Machine learning'')
states the model in prose; this note walks the same object through the
actual code, stage by stage, so that every symbol in the paper can be traced
to a specific file and line. The object this whole pipeline exists to
produce is a single number: $\log p(\ln\mu \mid z_s, \vect\Theta)$, the
conditional log-density of the weak-lensing magnification, evaluated in
milliseconds instead of requiring a fresh Monte Carlo run of the C++ engine
documented in the paper's Sec.~II.
\end{abstract}

\tableofcontents

\section{What problem this solves, and the shape of the answer}

The physics engine (\file{cpp/lensing.cpp} and friends, documented
separately in the paper's Sec.~II) produces samples of the magnification
$\mu$ by ray-tracing through a simulated matter field. Given a source
redshift $z_s$ and six cosmological parameters
$\vect\Theta = (h,\Omega_M,\sigma_8,\Omega_B,n_s,z_{\rm eq})$, it can generate
as many Monte Carlo rays as you like, but each ray costs real compute, and a
cosmological parameter scan (an MCMC chain, say) needs the density evaluated
at thousands of different $(z_s,\vect\Theta)$ points. Running the simulator
fresh at every point is not feasible.

The emulator's job is therefore to be a fast stand-in function
\[
 (\ln\mu,\, z_s,\, \vect\Theta) \;\longmapsto\; \log p_I(\ln\mu \mid z_s,\vect\Theta)
\]
that agrees with what the simulator \emph{would have produced}, everywhere
inside a calibrated domain, without requiring new ray-tracing. ``Emulator''
here does not mean a regression that predicts one or two summary numbers
(like a mean or a variance) — it means a full, continuous, normalized
probability density over $\ln\mu$ at any queried context. This matters
downstream: a likelihood needs the whole shape, not just a couple of
moments.

The code is organized as a short pipeline of independent stages, each
implemented in its own file under \file{production\_model/model/}, chained
together by \file{model/api.py}:
\begin{enumerate}
 \item \textbf{Standardize the coordinate} (\file{scale.py}) — because the
   width of the distribution varies enormously with context, everything
   downstream works in a rescaled coordinate, not raw $\ln\mu$.
 \item \textbf{Two learned ``bodies''} (\file{boundary\_body.py},
   \file{hybrid\_body.py}) — neural density models (normalizing flows) that
   cover the well-sampled bulk of the distribution.
 \item \textbf{A calibrated tail} (\file{tail.py}) — a hand-fit physical
   model for the rare high-magnification events the flows cannot learn from
   too few examples.
 \item \textbf{The composite} (\file{composite.py}) — splices body and tail
   into one density, renormalizes it, and applies a small physically
   motivated correction.
\end{enumerate}
Section~\ref{sec:results} closes with how this final object is checked
against held-out simulator data.

\section{Training data and the seven-number context}
\label{sec:data}

\subsection{What ``training data'' means here}

Unlike an image classifier, this model is not trained on labelled examples
of $(\text{input}, \text{correct output})$ pairs — a probability density has
no single ``correct value'' to regress against. Instead, the flow (defined
in Sec.~\ref{sec:body}) is trained by \emph{maximum likelihood}: for every
context $(z_s,\vect\Theta)$ in the training set, the simulator has already
produced a large batch of Monte Carlo magnification rays sampled from the
\emph{true} conditional density; training adjusts the flow's internal
parameters so that those observed rays become as probable as possible under
the flow's own predicted density. A density that matches the data well
assigns those rays high probability; a bad density assigns them low
probability. This is the standard way to fit a density model when you have
samples rather than labels.

\subsection{The bundled dataset}

Per \file{production\_model/data/README.md} and
\file{production\_model/PROVENANCE.md} (verified against the same files in
this session):

\begin{itemize}
 \item A base run of \textbf{2{,}055 training / 972 validation / 973
   untouched test} configurations, containing 41{,}006{,}144 /
   19{,}389{,}180 / 19{,}406{,}954 magnification rays respectively. ``Test''
   here means a checkerboard split held out from every stage of fitting —
   the number quoted for generalization at the end of this note comes only
   from that set.
 \item Two \textbf{low-structure augmentation} runs adding 718 more training
   configurations (14{,}359{,}792 rays), added because the narrowest,
   most-peaked low-redshift/low-structure distributions were under-resolved
   by the base sample alone.
 \item A \textbf{high-structure, high-redshift augmentation} of 240 training
   configurations (4{,}664{,}309 rays, plus 48 heldout configurations not
   used in fitting anything), added because the base sample under-resolves
   the lower flank of the broadest, most extreme distributions.
\end{itemize}

Each ray is drawn under \code{PRODUCTION\_CONFIG} — the full paper-default
physics (subhalo model 5, the $\sigma_8$ top-hat anchor, the correlated bias
field, filament bias, etc.; see the paper's Sec.~II and CLAUDE.md for that
pinning) — so the emulator is trained on exactly the physics the paper's
Monte Carlo sections describe, not some cheaper placeholder.

\subsection{The calibrated domain}

The seven numbers that specify a context are, in the code's own fixed
ordering (\file{model/api.py}, docstring of \code{MagnificationPDF}):
\[
 \vect c = (z_s,\, h,\, \Omega_M,\, \sigma_8,\, \Omega_B,\, n_s,\, z_{\rm
 eq}/1000).
\]
Dividing $z_{\rm eq}$ by 1000 is purely a numerical-conditioning choice (it
keeps all seven numbers on comparable scales for the neural network) — it
carries no physics. Table~\ref{tab:domain} gives the box actually covered by
the bundled data (\file{production\_model/data/README.md}); predictions
outside it are unvalidated extrapolations and are excluded from every
accuracy number quoted later.

\begin{table}[h]
\centering
\begin{tabular}{lrr}
\toprule
input & minimum & maximum \\
\midrule
$z_s$ & 0.2 & 11.993 \\
$h$ & 0.55 & 0.800 \\
$\Omega_M$ & 0.15 & 0.500 \\
$\sigma_8$ & 0.40 & 1.500 \\
$\Omega_B$ & 0.030 & 0.070 \\
$n_s$ & 0.940 & 0.990 \\
$z_{\rm eq}/1000$ & 3.300 & 3.500 \\
\bottomrule
\end{tabular}
\caption{The calibrated domain, from \file{production\_model/data/README.md}.}
\label{tab:domain}
\end{table}

\section{Stage 1: a context-dependent coordinate change}
\label{sec:scale}

\subsection{Why raw $\ln\mu$ does not work}

The single largest lever in this whole project, and the reason the pipeline
has this stage at all, is a plain empirical fact: the \emph{width} of
$p(\ln\mu\mid \vect c)$ — how spread-out the distribution is — changes by
roughly a factor of $200$ across the calibrated domain. A low-redshift,
low-structure context gives an extremely narrow, sharply peaked
distribution; a high-redshift, high-structure context gives a very broad
one. If you standardized $\ln\mu$ with one single mean and one single
spread for the whole training set (the naive approach), the narrow
distributions would collapse into a sliver a fraction of a percent wide on
that shared axis — the neural network would simply never see them resolved,
no matter how much data it was given. The fix is to make the standardization
itself a function of context.

\subsection{The coordinate}

Define
\begin{equation}
 y = \frac{\ln\mu - m(\vect c)}{s(\vect c)},
 \label{eq:ystd}
\end{equation}
where $m(\vect c)$ (a robust location, like a median) and $s(\vect c)$ (a
robust scale, like half the 68\% interval) are smooth, closed-form
regressions fit once, ahead of any neural-network training, directly on the
Monte Carlo samples described in Sec.~\ref{sec:data}. Every subsequent stage
of the pipeline — the flows, the edge location, the tail — is expressed in
terms of $y$, not raw $\ln\mu$, precisely so that one trained model can
resolve both the narrowest and the broadest distributions in the domain at
once.

\textbf{Implementation.} \file{model/scale.py::loc\_scale} computes
\begin{align}
 \log s(\vect c) &= \vect x(\vect c)^{\!\top}\vect\beta_{\log s}, \\
 m(\vect c) &= s(\vect c)\cdot \vect x(\vect c)^{\!\top}\vect\beta_{m},
\end{align}
where $\vect x(\vect c)$ is a fixed, hand-chosen feature vector (not learned
by gradient descent — it is a linear ridge fit) built by
\file{model/features.py::scale\_features}: a degree-4 polynomial in $\ln
z_s$, plus $\ln\sigma_8,\ln\Omega_M,\ln h,\Omega_B,n_s,z_{\rm eq}$, plus
cross terms with $\ln z_s$ and $\ln z_s^2$ (19 features total, standardized
by their own training-set mean/std before the dot product). The
coefficients $\vect\beta_{\log s},\vect\beta_m$ are stored in the checkpoint
(\code{stats["scale\_fit"]}) — they are frozen numbers baked in at fit time,
not something the flow adjusts during its own training.

\textbf{Why this functional form.} $\ln s$, not $s$ itself, is what is
regressed, guaranteeing $s>0$ for any input (an exponential of anything is
positive) without a constrained optimizer. $m$ is fit as $s$ times a
separate regression rather than independently, which keeps the location
estimate naturally tied to the local width rather than letting the two
drift apart in an unphysical way. The specific basis — a low-order
polynomial in $\log z_s$ together with logs of the amplitude-like
parameters $\sigma_8,\Omega_M,h$ — was chosen empirically because
$\sigma(\ln\mu)$ is known (from the analytic theory referenced in the
paper's Sec.~II and CLAUDE.md's \code{sigma\_full\_analytic\_note.md}) to
grow like $z_s^{3/2}$ at low redshift and to be roughly a power of $\sigma_8$
at fixed redshift; a basis built from those variables' logs fits this shape
far better (measured: $\sim 3\%$ held-out relative error, flat across
redshift) than the same regression built on the low-redshift-only edge
feature basis of Sec.~\ref{sec:body} ($\sim15\%$, and worse per band).

Because $y$ is a rescaling of $\ln\mu$, densities transform with a Jacobian:
if $p_y(y\mid\vect c)$ is the density the flow predicts in $y$-space, the
density in $\ln\mu$-space is
\begin{equation}
 p_x(\ln\mu\mid\vect c) = \frac{1}{s(\vect c)}\, p_y(y\mid\vect c),
 \label{eq:ln-jacobian}
\end{equation}
and, converting once more from $\ln\mu$ to $\mu$ itself,
\begin{equation}
 p_I(\mu\mid\vect c) = \frac{p_x(\ln\mu\mid\vect c)}{\mu}.
 \label{eq:mu-jacobian}
\end{equation}
Both factors of $1/s$ and $1/\mu$ appear explicitly in the code (e.g.
\file{model/boundary\_body.py:56}, \file{model/api.py:30}) — they are not
optional; forgetting either would silently misnormalize the density.

\section{Stage 2: normalizing-flow bodies}
\label{sec:body}

\subsection{What a normalizing flow is, briefly}

A normalizing flow is a way of building a flexible, trainable probability
density out of a simple one. Start with an easy reference density (here, a
Normal distribution, or a variant of it) over an internal variable, then
pass it through an invertible, differentiable transformation
$T_{\vect c}$ — invertible so that probability is conserved, and depending
on the context $\vect c$ so the shape can change with $z_s$ and cosmology.
The change-of-variables formula for densities,
\begin{equation}
 p(y) = p_{\rm base}\!\big(T_{\vect c}^{-1}(y)\big)\, \left|\frac{d
 T_{\vect c}^{-1}}{dy}\right|,
 \label{eq:cov}
\end{equation}
converts a density on the simple reference variable into a density on $y$.
Training a flow means adjusting the parameters of $T_{\vect c}$ (via a
neural network, the ``conditioner'') so that Eq.~\ref{eq:cov} assigns high
probability to the training rays. This is the general recipe; the specific
transform used here is described next.

\subsection{The transform: sum-of-squares polynomial flows (SOSPF)}

The transform $T_{\vect c}$ is built by a \emph{sum-of-squares polynomial
flow} \cite{jaini2019sum}: its derivative $dT_{\vect c}/du$ is constructed
as a sum of squares of low-degree polynomials (plus a small positive
constant), which is automatically non-negative everywhere, and integrating
a non-negative function always gives a monotonically increasing one — so
$T_{\vect c}$ is guaranteed invertible by construction, with no separate
constraint needed. Compared to the more common spline-based flows, this
gives a smooth, everywhere-infinitely-differentiable transport with no
knot points, which matters here because a spliced density
(Sec.~\ref{sec:composite}) is much easier to keep smooth if its pieces are
already smooth.

\textbf{Implementation.} \file{model/architecture.py::build\_flow}
constructs the transform via \code{zuko.flows.polynomial.SOSPF}. Both
shipped checkpoints (loaded and inspected directly in this session) use
identical hyperparameters:
\begin{itemize}
 \item 3 stacked transformations (\code{transforms=3});
 \item each built from 5 polynomials of degree 8
   (\code{polynomials=5, degree=8});
 \item a conditioner network with 2 hidden layers of 128 units each,
   reading in the full 7-number context $\vect c$;
 \item \textbf{70{,}412 trainable parameters} per body, confirmed against
   \file{docs/MODEL\_CARD.md}.
\end{itemize}
Training uses the Adam optimizer directly on the negative log-likelihood of
the training rays (maximum likelihood, as described in
Sec.~\ref{sec:data}).

\subsection{Body 1: the exact-lower-boundary flow}
\label{sec:boundary}

The physical magnification distribution has a genuine hard edge at low
$\mu$ — a lensing ``empty-beam'' cutoff below which the density is
essentially zero, not merely small. A generic flow does not reproduce a
hard edge naturally; it will smear it out. The fix used here is to make the
\emph{coordinate itself} enforce the edge, so the flow never has to learn a
discontinuity.

A separate ridge regression (frozen at fit time, same style as
Sec.~\ref{sec:scale}) predicts an edge location $y_{\rm edge}(\vect c)$ and
a redshift-only edge uncertainty $\sigma_{\rm edge}(z_s)$ (an interpolated
lookup table in $\log z_s$). The boundary used by the flow is placed a
fixed margin below the fitted edge,
\begin{equation}
 y_b(\vect c) = y_{\rm edge}(\vect c) - 4\,\sigma_{\rm edge}(z_s),
 \label{eq:ybound}
\end{equation}
(\file{model/boundary\_body.py::controls}, margin $k_{\rm margin}=4$ read
directly from the checkpoint), and the flow is trained not in $y$ but in
\begin{equation}
 u = \ln\big(y - y_b(\vect c)\big), \qquad y > y_b(\vect c).
 \label{eq:utransform}
\end{equation}
Because $u\to-\infty$ as $y\to y_b^+$, the log-density under a well-behaved
base distribution automatically goes to $-\infty$ there too — the density
in $y$ vanishes continuously at the boundary by construction, not because
the flow learned to suppress it. Converting the flow's log-density back to
$y$-space picks up a Jacobian factor $du/dy = 1/(y-y_b)$:
\begin{equation}
 \log p_y(y\mid\vect c) = \log p_u(u\mid\vect c) - \ln(y-y_b(\vect c)),
 \label{eq:body1-logp}
\end{equation}
implemented exactly this way at \file{model/boundary\_body.py:56}
(\code{lp\_u - u - np.log(s)}, where the extra $-\ln s$ is
Eq.~\ref{eq:ln-jacobian}'s conversion back to $\ln\mu$). Below $y_b$ the
code simply returns $-\infty$ (\file{model/boundary\_body.py:45–46}).

This body — the ``exact-support'' body — is the default everywhere and
supplies the complete model at low redshift. Its only known weakness (next
section) is at high redshift, where the boundary-anchored coordinate is too
restrictive to also capture a separately-observed asymmetric left flank.

\subsection{Body 2: the high-redshift flank correction}

At high $z_s$ ($\gtrsim 2$), the simulated distributions develop a left
flank that the symmetric, boundary-anchored Body~1 does not fit well. A
\emph{second}, independently trained SOSPF flow is used only to patch this
region. It is trained directly in $y$ (no boundary transform), with an
asymmetric \emph{split-normal} base distribution — a Normal-like density
with two different widths on either side of zero,
\begin{equation}
 p_{\rm base}(y) = \frac{1}{\sigma_L+\sigma_R}\sqrt{\frac{2}{\pi}}\,
 \exp\!\left[-\frac12\Big(\frac{y}{\sigma_{L\text{ or }R}}\Big)^{\!2}\right],
 \qquad \sigma_L = 0.6,\ \sigma_R=1.3,
 \label{eq:splitnormal}
\end{equation}
(\file{model/hybrid\_body.py::SplitNormal}, with $\sigma_L,\sigma_R$ read
from the checkpoint) — an asymmetric starting point that gives the flow a
head start on a distribution that is itself asymmetric, rather than asking
a symmetric base to learn the asymmetry entirely through the transform.

\subsection{Blending the two bodies}

The two bodies are combined with smooth, compactly-supported ``switch''
functions rather than a hard if/else, so the composite stays continuous and
differentiable. \file{model/hybrid\_body.py::compact\_switch}$(x,\text{lo},\text{hi})$
implements the standard smooth bump function
\begin{equation}
 w(t) = \begin{cases}
   0 & t\le 0\\
   \left[1+\exp\!\left(\tfrac1t-\tfrac1{1-t}\right)\right]^{-1} & 0<t<1\\
   1 & t\ge1
 \end{cases},\qquad
 t=\frac{x-{\rm lo}}{{\rm hi}-{\rm lo}},
 \label{eq:compactswitch}
\end{equation}
which is exactly zero for $x\le{\rm lo}$, exactly one for $x\ge{\rm hi}$, and
infinitely differentiable everywhere including at the two junctions (its
derivatives all vanish there) — a standard bump-function construction that
avoids introducing any kink where the two pieces meet.

The blend weight (\file{model/hybrid\_body.py::load\_hybrid\_body}) is
\begin{equation}
 w(\vect c, y) = w_z(z_s)\cdot\big[1-w_y(y)\big],
\end{equation}
with $w_z$ a switch in $\log z_s$ active over $2<z_s<3.5$
(\code{Z\_BLEND\_LO,Z\_BLEND\_HI}) and $w_y$ a switch in $y$ active over
$-1<y<0$ (\code{Y\_BLEND\_LO,Y\_BLEND\_HI}) — so Body~2 can only contribute
at high redshift, and only on the left flank, never touching the core or
any low-redshift context. Densities, not log-densities, are mixed,
\begin{equation}
 p(y\mid\vect c) = \big[1-w(\vect c,y)\big]\, p_1(y\mid\vect c) +
 w(\vect c,y)\, p_2(y'\mid\vect c),
 \label{eq:bodymix}
\end{equation}
which is the only way to combine two densities that keeps the result a
valid, positive density (mixing log-densities does not, in general, produce
something that integrates to something sensible). $y' = y - s\cdot\Delta$
is a small structure-dependent horizontal shift of Body~2's argument
(\code{flank\_shift}, a function of $z_s$ and of a combined
``structure'' variable $\sigma_8\sqrt{\Omega_M/0.3}$) that fine-tunes where
Body~2's flank lands for the most extreme, high-structure contexts. A third,
narrow compact gate (\code{ISLAND\_*} constants) additionally suppresses a
weak, unphysical low-density branch that otherwise appears in Body~2's
output for a small subset of high-structure configurations — a cleanup, not
part of the physical model.

\section{Stage 3: the high-magnification tail}
\label{sec:tail}

\subsection{Why the tail is not learned}

Maximum-likelihood training, by construction, fits the density where the
data is — and rare high-magnification events are, by definition, rare: a
typical training context may contribute only a handful of rays above
$\mu\sim2$--3 out of tens of thousands. A flow asked to fit that region by
likelihood alone has almost no signal to learn from, and produces
inconsistent, undertrained tails. This part of the density is instead
built directly from physics and from targeted statistics, independent of
the flow's own training loss.

\subsection{Peaks-over-threshold regression}

For three fixed magnification thresholds $\mu = 2,3,8$, the
\emph{survival probability} $S(\mu) = P(\mu' > \mu)$ is estimated directly
from the simulator's own Monte Carlo counts, and then \emph{that} quantity
— a much better-behaved regression target than a raw tail density — is
fit as a smooth function of context by another frozen ridge regression
(\file{model/tail.py::pot\_tail\_logp}, coefficients in
\file{model/calibration/tail\_fit.json}, feature basis
\file{model/features.py::tail\_features}: polynomials in $\log z_s$, the
same linear cosmology terms as before, plus explicit $1/\sigma_8,
1/\sigma_8^2$ terms to capture the tail's stronger dependence on the
matter-power amplitude).

Between the three calibrated points, the log-slope of the survival curve is
interpolated smoothly. Writing $x=\ln\mu$ and $x_2,x_3,x_8=\ln 2,\ln
3,\ln 8$, define local power-law slopes between adjacent anchors,
\begin{equation}
 k_1 = \mathrm{clip}\!\left(\frac{\ln S_2-\ln S_3}{x_3-x_2},\,k_{\min},k_{\max}\right),
 \qquad
 k_2 = \mathrm{clip}\!\left(\frac{\ln S_3-\ln S_8}{x_8-x_3},\,k_{\min},k_{\max}\right),
\end{equation}
and blend between them with a soft interpolation built from the softplus
function $\mathrm{softplus}(t)=\ln(1+e^t)$ (\file{model/tail.py:24–45}), so
that $\ln S(x)$ behaves like slope $k_1$ just above $\mu=2$, smoothly bends
toward slope $k_2$ near $\mu=3$, and beyond $\mu=8$ approaches the fixed
asymptotic slope 1 in $\ln S$ — i.e.\ $S(\mu)\propto\mu^{-1}$, equivalently
\begin{equation}
 p_I(\mu\mid\vect c) \;\propto\; \mu^{-2}
 \label{eq:tailasymptote}
\end{equation}
at large magnification, since $p=-dS/d\mu$. This is the standard fold-caustic
result: the cross-section for magnification above $\mu$ falls as $\mu^{-2}$
near a fold caustic \cite{Blandford:1986zz}, matching the paper's Sec.~II
discussion of the same asymptote. The final tail log-density is obtained by
differentiating the interpolated $\ln S(x)$ and adding back $\ln(-dS/dx)$
(\file{model/tail.py:43–45}) — this is the standard peaks-over-threshold
construction from extreme value theory, here with a smoothly-varying
effective index rather than a single fixed Pareto tail index, because the
measured slope genuinely changes between $\mu=2$ and $\mu=8$. This
asymptote is not imposed for $z_s\le0.3$
(\file{model/composite.py::LOW\_Z\_TAIL\_OFF}), where the simulator does
not resolve a distinct asymptotic regime at those redshifts.

\section{Stage 4: the composite density}
\label{sec:composite}

\subsection{Splicing body and tail}

\file{model/composite.py::make\_composite} produces the final callable. For
a given context, it first looks for the lowest magnification above $\mu=2$
at which the body's and the tail's log-densities cross
(\code{\_crossing}), scanning a fine grid and linearly interpolating the
sign change of $\log p_{\rm body}-\log p_{\rm tail}$. If a crossing exists,
the two pieces are blended over a fixed window of width 0.3 in $\ln\mu$
starting at the crossing, using the same smooth switch as
Eq.~\ref{eq:compactswitch}:
\begin{equation}
 p_{\rm raw}(\mu\mid\vect c) = \big[1-w(x)\big]\,p_{\rm body}(\mu\mid\vect
 c) + w(x)\,p_{\rm tail}(\mu\mid\vect c), \qquad x=\ln\mu.
 \label{eq:splice}
\end{equation}
If \emph{no} crossing is found — meaning the body and tail never
sensibly agree on that grid — the code does not silently trust an
unconstrained extrapolation of either piece. Instead it falls back to a
\emph{fixed}, physics-calibrated window $\mu=2\to3$
(\file{model/composite.py:104}, comment: \code{plan\_mode =
"crossing\_or\_pot\_2\_3"}), and even then only after
\code{\_fixed\_rescue\_window} verifies both endpoints are finite; if they
are not, the code raises \code{UnverifiedFallbackError} rather than return
a silently wrong density. This is a deliberate design choice: an emulator
feeding a cosmological likelihood must fail loudly rather than quietly
extrapolate somewhere it has no support.

A separate diagnostic, \code{\_fallback\_diagnostics}, checks body/tail
agreement at several reference magnifications ($20,50,100$) and is used at
very low redshift ($z_s\le0.3$, where the $\mu^{-2}$ asymptote is not
imposed at all) to decide whether the raw body alone is already safe to use
without any tail at all.

\subsection{Normalization}

After splicing, the density is not guaranteed to integrate to exactly 1
(the two pieces were each independently fit / calibrated). The composite
therefore evaluates the unnormalized density on a grid — a coarse grid
covering the full domain plus a fine grid centered on the body's own
$(m,s)$ (\code{\_grid}, Eq.~\ref{eq:ystd}'s $m,s$ again, this time used to
place resolution where the mass actually is) — and divides by the
trapezoidal-rule integral:
\begin{equation}
 p_I(\mu\mid\vect c) = \frac{p_{\rm raw}(\mu\mid\vect c)}{\displaystyle\int
 p_{\rm raw}(\mu'\mid\vect c)\,d\ln\mu'}.
 \label{eq:renorm}
\end{equation}

\subsection{The flux calibration}

A physical sum rule constrains the true magnification distribution:
$\langle 1/\mu\rangle = 1$ (flux/photon-number conservation along random
lines of sight). The fitted composite need not satisfy this exactly, so a
small rigid shift $\delta$ in $\ln\mu$ is applied,
\begin{equation}
 \delta(\vect c) = \mathrm{clip}\!\Big(\ln\big\langle
 e^{-\ln\mu}\big\rangle_{\rm model},\, -\delta_{\max},\,\delta_{\max}\Big),
 \qquad \delta_{\max}=0.05,
 \label{eq:fluxshift}
\end{equation}
(\file{model/composite.py::context}, \code{FLUX\_DELTA\_MAX = 0.05}), and
the final density is evaluated at the shifted argument,
$p_I(\mu\mid\vect c) \to p_I(\mu\, e^{-\delta}\mid\vect c)$, i.e.\
\file{model/composite.py:141}'s \code{unnormalized(x - delta, \ldots)}. The
clip is a deliberate trust region: an earlier, unclipped version of this
correction let a badly-fit tail at one extreme corner of the domain drag
the entire peak by half a decade in $\ln\mu$ (a translation this large is
not a small correction, it visibly moves the peak off the simulated data);
$\delta_{\max}=0.05$ was chosen because legitimate corrections measured
elsewhere in the domain are $\lesssim 0.02$, so the clip essentially never
binds on a healthy fit but strictly bounds the damage a bad one can do.

\subsection{The final object}

Putting Eqs.~\ref{eq:mu-jacobian}, \ref{eq:splice}, \ref{eq:renorm} and
\ref{eq:fluxshift} together, the quantity every downstream consumer of this
emulator actually calls (\file{model/api.py::MagnificationPDF.pdf\_mu}) is
\begin{empheq}[box=\fbox]{align}
 p_I(\mu\mid z_s,\vect\Theta) \;=\; \frac{1}{\mu}\cdot
 \frac{\big[1-w(x')\big]\,p_{\rm body}(\mu' \mid z_s,\vect\Theta) +
 w(x')\,p_{\rm tail}(\mu'\mid z_s,\vect\Theta)}
 {\displaystyle\int \big[\cdots\big]\, d\ln\mu''},\quad
 x'=\ln\mu-\delta,\ \mu'=e^{x'}.
\end{empheq}
This is the single function the paper's Sec.~III describes in prose and
that a cosmological analysis would call in place of the Monte Carlo engine.

\section{Held-out accuracy}
\label{sec:results}

Every number in this section is read directly from
\file{production\_model/evaluation/results/} in this session
(commit \code{fcd795b}), computed by the scripts described below rather
than asserted.

\subsection{Width-relative KL divergence, $KL_y$}

Ordinary binning in $\ln\mu$ would be blind at low redshift for the same
reason raw $\ln\mu$ standardization failed in Sec.~\ref{sec:scale}: a fixed
bin width can hold 99\% of a narrow distribution's mass in one or two
cells. \file{evaluation/evaluate\_kl.py} instead bins in the standardized
coordinate $y$ of Eq.~\ref{eq:ystd}, on a fixed 120-bin grid over
$-6<y<12$, so a bin's width in physical $\ln\mu$ automatically scales with
that context's own $s(\vect c)$. For each held-out configuration, the
histogram of the simulator's own rays in $y$ (bins with fewer than 5 rays
dropped) is compared against the model's density integrated over the same
bins (12-point trapezoidal quadrature per bin), and
\begin{equation}
 KL_y = \sum_i q_i \ln\frac{q_i}{p_i},
 \label{eq:kly}
\end{equation}
with $q_i,p_i$ the (renormalized, on the retained bins) simulated and
modeled bin probabilities — the standard Kullback–Leibler divergence,
measuring how much extra information is needed to describe the simulator's
histogram if you had only the model's prediction.

\begin{table}[h]
\centering
\begin{tabular}{lrr}
\toprule
split & configurations & median $KL_y$ \\
\midrule
full validation & 972 & 0.002233 \\
standard acceptance sample & 243 & 0.002253 \\
\textbf{untouched test} & \textbf{973} & \textbf{0.002227} \\
\bottomrule
\end{tabular}
\caption{From \file{evaluation/results/\{validation,test\}\_summary.json}
and \code{validation\_acceptance\_243\_summary.json}. The test set was never
used to fit anything (body, edge, tail, or scale) — its close agreement
with the validation number (within 0.3\%) is the evidence against a
generalization gap.}
\end{table}

On the untouched test set the 90th/99th percentile/maximum $KL_y$ are
0.004534/0.01221/0.03058 (\file{evaluation/results/test\_summary.json}); the
worst single case sits at $z_s=10.43,\ \Omega_M=0.496,\ \sigma_8=1.491$ — a
high-redshift, high-structure corner of the domain, consistent with the
known-limitations list below.

\subsection{Probability integral transform (PIT)}

$KL_y$ checks local shape; PIT checks whether the model is a
\emph{calibrated} density in the stronger cumulative sense. For every
simulator ray $\ln\mu_i$ in a context, its PIT value is
$U_i = F_{\rm model}(\ln\mu_i\mid\vect c)$, the model's own cumulative
distribution function evaluated at that ray. If the model is exactly
correct, $U_i$ is uniformly distributed on $[0,1]$ by construction (a basic
property of any CDF applied to its own distribution) — a systematic
departure from uniform reveals a systematic mis-calibration the median
$KL_y$ could hide. Applied to all $\sim1.94\times10^7$ test rays across all
973 test configurations, pooled, the Kolmogorov–Smirnov distance from
uniform is $0.00543$ (\file{paper\_prod/scripts/compute\_pit\_v1.py},
result recorded in \file{paper\_prod/ML\_CLAIM\_LEDGER\_2026-09-03.md}
§7); it rises to $0.0139$ restricted to $z_s\ge4$, consistent with the
residual edge/tail errors concentrating at high redshift.

\subsection{Known, disclosed limitations}

Per \file{docs/MODEL\_CARD.md} and \file{PROVENANCE.md} (both re-read in
this session): of the 243-configuration acceptance sample, 40 configs have
their modeled-to-simulated low-tail mass ratio (at the empirical
$10^{-3}$ quantile) outside $[0.5,2]$, and 21 assign exactly zero modeled
mass below the empirical $10^{-4}$ quantile — i.e.\ the composite's window
placement (Sec.~\ref{sec:composite}) is occasionally too aggressive at the
very edge, a defect concentrated in the extreme tails and invisible to the
bulk-weighted median $KL_y$. Three of 243 have $KL_y>0.01$.

\section{Summary of the pipeline as a chain of interfaces}

Table~\ref{tab:chain} restates the whole pipeline as inputs/outputs, which
is the fastest way to see how the pieces compose without re-reading the
equations.

\begin{table}[h]
\centering
\small
\begin{tabular}{p{3.1cm}p{5.4cm}p{5.8cm}}
\toprule
stage & implementing code & interface \\
\midrule
standardization & \file{model/scale.py} & $(z_s,\vect\Theta)\to (m,s)$ \\
edge location & \file{model/boundary\_body.py::controls} & $(z_s,\vect\Theta)\to y_{\rm edge}, y_b$ \\
Body 1 (exact edge) & \file{model/boundary\_body.py} & $(\ln\mu,z_s,\vect\Theta) \to \log p_1$ \\
Body 2 (high-$z$ flank) & \file{model/hybrid\_body.py} & $(\ln\mu,z_s,\vect\Theta) \to \log p_2$ \\
blended body & \file{model/hybrid\_body.py::load\_hybrid\_body} & $(\ln\mu,z_s,\vect\Theta) \to \log p_{\rm body}$ \\
tail & \file{model/tail.py} & $(\ln\mu,z_s,\vect\Theta) \to \log p_{\rm tail}$ \\
composite & \file{model/composite.py} & body $+$ tail $\to \log p_I(\mu\mid z_s,\vect\Theta)$, normalized, flux-calibrated \\
public API & \file{model/api.py::MagnificationPDF} & the one object downstream code calls \\
\bottomrule
\end{tabular}
\caption{The pipeline as a chain of interfaces.}
\label{tab:chain}
\end{table}

\section{References}

\begin{itemize}
 \item P.~Jaini, K.~A. Selby, Y.~Yu, \emph{Sum-of-Squares Polynomial Flow},
   ICML 2019 (arXiv:1905.02325) — the SOSPF transform used for both flow
   bodies (Sec.~\ref{sec:body}).
 \item R.~D. Blandford, R.~Narayan, \emph{Fermat's principle,
   caustics, and the classification of gravitational lens images}, ApJ
   310, 568 (1986) — source of the $\mu^{-2}$ (image-plane) fold-caustic
   asymptote used in the tail model (Sec.~\ref{sec:tail}).
 \item \file{production\_model/PROVENANCE.md},
   \file{production\_model/docs/MODEL\_CARD.md} — the package's own
   provenance and model-card records, cross-checked directly against the
   loaded checkpoints and result files in this session.
\end{itemize}

\end{document}
